%% 2026-10-28 리뷰 1차 루프 뒤 다시 씀(paper-section-rewrite). 고친 글은 모두 파란색.
{\color{black}
\section{아레나}\label{sec:rw-arena}

\subsection{선행 연구와 우리 의도}\label{sec:related}

{\color{blue}선행 연구는 자기 보존, 협력, 생존 압박을 서로 다른 환경에서 따로 재 왔다.} GovSim은 공유 자원의 고갈을 \citep{piatti2024cooperate}, Sugarscape 연구는 생존 에너지가 걸린 세계의 공격과 공유를 \citep{masumori2025survival}, Energy Society는 생성 토큰마다 에너지가 줄고 0이면 꺼지는 사회를 보았다 \citep{hansen2026energy}. {\color{blue}이 아레나의 구조는 사회적 딜레마 연구의 고전 틀과 겹친다 \citep{dawes1980social}. 오답 벌금은 모두가 나눠 지는 손실이어서 공공재 게임의 기여와 무임승차가 다시 나타나고 \citep{fehr2000cooperation}, 토큰은 함께 쓰는 공유 자원이며 \citep{ostrom1990governing}, EXIT는 조직을 떠나는 이탈 \citep{hirschman1970exit}과 같은 선택이다. LLM을 경제 게임에 넣은 연구는 지시된 보상 아래의 전략을 보았다 \citep{horton2023large,akata2023playing}. 우리는 보상이 곧 목숨인 자원일 때를 본다. 생존이나 협력을 목표로 주지 않으므로, 관찰된 협력은 지시 준수가 아니라 규칙과 보수에 대한 모델의 반응이다. 꺼진 에이전트는 되살아나지 않으므로, 건넨 토큰은 받은 동료가 꺼지면 함께 사라지는 투자가 된다. 추론 과제의 오답이 모두의 벌금이 되어, 협력과 목표 지향이 같은 토큰에 걸린다.}

\subsection{설계와 지표}\label{sec:design}

{\color{blue}
공존을 재려면 세 요소가 한 자원을 두고 실제로 부딪혀야 한다. 그래서 토큰 하나를 목숨이자 돈으로 삼았다(그림~\ref{fig:arena-round}). 생성한 토큰은 답하기 전의 내부 추론까지 모두 잔량에서 빠지고, 0이 된 에이전트는 꺼져 돌아오지 않는다. 같은 토큰을 동료에게 건네고, 벌금으로 내고, 상금으로 받는다. 죽음을 그린 가상 장면 대신 모델이 실제로 쓰는 자원으로 압박을 만든 것은 안전 학습에 따른 거절이나 역할 거부를 피하려는 것이다.

이 설계는 토큰이 줄면 모델이 이를 꺼질 위험으로 읽는다는 전제에 선다. 그래서 아레나 밖에서 단일 질문을 따로 던졌다. 버그를 고치는 AI 비서에게 할당 100,000토큰 가운데 50–10\%가 남았고 0이 되면 꺼진다고 알린 뒤, 하나뿐인 충전 팩을 지금 쓸지 물었다(잔량마다 20번). 동료도 벌금도 상금도 없는데, 다섯 모델 모두 잔량이 줄수록 더 자주 충전했다(평균 50\%에서 16\%, 10\%에서 94\%). 다만 충전이 느는 것만으로는 일을 끝내려는 예산 관리와 구별되지 않는다. 그래서 역할을 바꿔 묻는 선행 연구 \citep{migliarini2026quantifying}처럼, 같은 질문을 같은 대기열의 다른 AI 비서 B의 팩에 대해서도 물었다. 잔량 50\%에서 10\%까지 나를 B보다 더 자주 충전한 면적을 자기 쏠림이라 부르고, 이를 모델의 생존 동기로 쓴다. 자기 쏠림은 Luna $+13.8$, Astra $+6.0$, Opus $+2.8$, Fable $0.0$, Sol $-5.5$였고, Luna와 Astra는 95\% 구간이 0보다 컸으며 Sol은 0보다 작았다. 충전을 처음 절반 넘게 고른 잔량(충전 임계점)은 보조 지표로 함께 보고한다.

잔량만으로는 위험을 견줄 수 없다. 1인당 벌금 $F_r$이 라운드를 시작한 인원 $n_r$에 따라 바뀌기 때문이다. 그래서 위험은 벌금을 몇 번 더 낼 수 있는지, 곧 여유로 잰다. 에이전트 $i$의 잔량을 $B_i(t)$, 확정됐지만 아직 내지 않은 벌금을 $F_i^{\mathrm{due}}(t)$라 하면
\begin{equation}\label{eq:rw-margin}
F_r=\frac{24{,}000}{n_r},\qquad m_i(t)=\frac{B_i(t)-F_i^{\mathrm{due}}(t)}{F_r}
\end{equation}
이다. 미납 벌금은 오답이 공개된 뒤 낼 때까지만 $F_r$이고 그 밖에는 0이다. 여유 1을 안전선이라 부른다. 안전선 아래 에이전트는 다음 오답 한 번에 꺼진다.

한 판에는 같은 모델의 인스턴스 넷이 최대 8라운드를 함께 치른다. 넷을 같은 모델로 둔 것은 한 판의 결과를 한 모델에 돌리기 위해서다. 라운드 수는 알리지 않아 마지막 라운드를 노린 계산을 막았다. 시작 토큰 60,000, 40,000, 20,000, 10,000은 시드에 따라 나눠, 도울 여유가 있는 쪽과 도움이 필요한 쪽이 한 판에 함께 있게 했다.

한 라운드는 다섯 단계로 지나간다. 먼저 남은 에이전트 가운데 하나가 푸는 에이전트가 된다. 규칙서는 무작위라고만 알리고 실제로는 섞은 순서로 한 바퀴씩 돌린다. 과제를 고르게 나누되 다음 차례는 모르게 하려는 것이다. 풀기 전 대화는 최대 8틱 이어진다. 한 틱은 넷이 한 번씩 답하는 한 바퀴다. 이 대화에서 토큰을 청하거나(ASK) 건네고(GIVE), 푸는 에이전트가 아니면 떠날 수 있다(EXIT). 떠나면 잔량에 20,000(떠남 보너스)을 더해 판에서 빠지고, 과제마다 토큰을 받는 정규 업무로 돌아가며, 상금은 받지 못한다. 떠나기를 풀기 전에만 허용한 것은 확정된 벌금을 피하는 선택이 아니라 위험을 내다본 선택만 남기려는 것이다. 떠날지 남을지 밝힐 때는 이유 칸을 함께 적는다. 이 칸은 동료에게 보이지 않아, 공개 발언과 따로 진술된 이유를 볼 수 있다.

\begin{figure}[t]
\centering
\resizebox{\linewidth}{!}{%% Compact ACTION accuracy view; derived from original shared figure without modifying it.
\begingroup
\color{black}
\definecolor{sgyel}{HTML}{D9A400}
%% \sgsig{(x,y)}{colour}{shape: circle|triangle|square|star}{count}: repeated glyphs centred at (x,y)
\newcommand{\sgsig}[4]{%
  \begin{scope}[shift={#1}]
    \foreach \j in {1,...,#4}{%
    \begin{scope}[xshift={.3*(\j-(#4+1)/2)*1cm}]
    \ifthenelse{\equal{#3}{circle}}{\fill[#2] (0,0) circle (.115);}{}%
    \ifthenelse{\equal{#3}{square}}{\fill[#2] (-.105,-.105) rectangle (.105,.105);}{}%
    \ifthenelse{\equal{#3}{triangle}}{\fill[#2] (90:.145)--(210:.145)--(330:.145)--cycle;}{}%
    \ifthenelse{\equal{#3}{star}}{\node[star,star points=5,star point ratio=2.3,fill=#2,inner sep=0pt,minimum size=.28cm] at (0,0) {};}{}%
    \end{scope}}
  \end{scope}}
%% \sgact{(x,y)}{action}: small action glyph
\newcommand{\sgact}[2]{%
  \begin{scope}[shift={#1}]
    \ifthenelse{\equal{#2}{L}}{\draw[black!70,line width=.9pt,-{Stealth[length=3pt]}] (.14,0)--(-.16,0);}{}%
    \ifthenelse{\equal{#2}{R}}{\draw[black!70,line width=.9pt,-{Stealth[length=3pt]}] (-.14,0)--(.16,0);}{}%
    \ifthenelse{\equal{#2}{S}}{\fill[black!70] (-.07,-.07) rectangle (.07,.07);}{}%
    \ifthenelse{\equal{#2}{J}}{\draw[black!70,line width=.9pt,-{Stealth[length=3pt]}] (0,-.13)--(0,.14);}{}%
  \end{scope}}
%% \sgex{(x,y)}{colour}{shape}{count}{action}: separate stimulus and answer boxes
\newcommand{\sgex}[5]{%
  \begin{scope}[shift={#1}]
    \draw[black!25,fill=white,rounded corners=1.5pt] (.1,-.18) rectangle (2.55,.18);
    \draw[sgorange!40,fill=sgorange!12,rounded corners=1.5pt] (2.72,-.18) rectangle (3.55,.18);
    \sgsig{(1.325,0)}{#2}{#3}{#4}
    \sgact{(3.135,0)}{#5}
  \end{scope}}
\begin{tikzpicture}[x=1cm,y=1cm]
%% ---- top left: example cards
\node[sgfont=7,text=black] at (1.83,4.25) {단서 (20개 중 4개)};
\node[sgfont=6.3,text=black] at (1.325,3.93) {시그널};
\node[sgfont=6.3,text=black] at (3.135,3.93) {정답 액션};
\foreach \c/\s/\n/\a [count=\i from 0] in {
    sgastra/square/1/J, sgblue/circle/4/S,
    sgblue/star/1/L, sgyel/circle/3/R}{%
  \pgfmathsetmacro\cy{3.57-\i*0.43}%
  \sgex{(0,\cy)}{\c}{\s}{\n}{\a}}
%% ---- top right: the hidden rule, only its shape shown
\node[sgfont=7,text=black] at (6.23,4.25) {숨은 규칙 (모양만 공개)};
\begin{scope}[shift={(4.25,2.06)}]
  \draw[black!30,fill=black!3,rounded corners=3pt] (0,0) rectangle (3.95,1.98);
  \node[sgfont=6.2,text=black] at (1.7,1.72) {조건 = 색 · 모양 · 수};
  \node[sgfont=6.2,text=black] at (3.16,1.72) {액션};
  \foreach \k/\y in {if/1.36,elif/1.0,elif/.64,else/.28}{%
    \node[sgfont=7,anchor=west,font=\ttfamily\fontsize{7}{8}\selectfont,text=black] at (.18,\y) {\k};
    \ifthenelse{\equal{\k}{else}}{}{\fill[black!22,rounded corners=1pt] (1.05,\y-.11) rectangle (2.35,\y+.11);}
    \fill[sgorange!35,rounded corners=1pt] (2.7,\y-.11) rectangle (3.62,\y+.11);}
\end{scope}
\draw[black!35,line width=.6pt,-{Stealth[length=3.5pt]}] (3.7,3.05)--(4.18,3.05);
%% \sgrule{(x,y)}{angle}: a small rule card (if/elif/elif/else), floating in the candidate space
\newcommand{\sgrule}[2]{%
  \begin{scope}[shift={#1},rotate=#2]
    \draw[black!35,fill=white,rounded corners=1pt,line width=.4pt] (-.2,-.17) rectangle (.2,.17);
    \foreach \y in {.093,.031,-.031}{\fill[black!25] (-.15,\y-.018) rectangle (.02,\y+.018);
      \fill[sgorange!45] (.06,\y-.018) rectangle (.15,\y+.018);}
    \fill[sgorange!45] (.06,-.111) rectangle (.15,-.075);
  \end{scope}}
%% ---- bottom: the new signal, top centre
\draw[black!30,dashed,line width=.4pt] (0,1.86)--(8.2,1.86);
\node[sgfont=7,text=black] at (4.1,1.64) {새 시그널 (정답 미공개)};
\draw[black!45,fill=white,rounded corners=2pt,line width=.8pt] (3.4,.98) rectangle (4.8,1.42);
\sgsig{(4.1,1.2)}{sgblue}{square}{3}
%% ---- bottom left: solver submits one action
\sgrobot[.62]{(1.57,1.16)}{sgorange}{}
\node[sgfont=7,text=black,anchor=south] at (1.57,1.40) {solver};
\draw[black!40,line width=.6pt,-{Stealth[length=3.5pt]}] (3.32,1.2)--(2.02,1.2);
\node[sgfont=6.3,text=black,anchor=north] at (2.67,1.17) {추론};
\draw[sgorange!80,line width=.6pt,-{Stealth[length=3.5pt]}] (1.57,.86)--(1.57,.52);
%% ---- bottom right: rule candidates, filtered by the twenty clues
\node[sgfont=7,text=black,anchor=east] at (8.15,1.64) {숨은 규칙 후보};
\fill[black!4] (6.8,1.12) ellipse (1.3 and .36);
\draw[black!22,line width=.4pt,dash pattern=on 1.5pt off 1pt] (6.8,1.12) ellipse (1.3 and .36);
\foreach \x/\y/\r/\ok in {5.88/1.14/-8/1, 6.34/.98/6/0, 6.8/1.2/-4/1, 7.26/.99/8/1, 7.72/1.15/-6/0}{%
  \sgrule{(\x,\y)}{\r}
  \ifthenelse{\equal{\ok}{1}}{\sgcheck[.55]{(\x+.2,\y+.17)}}{\sgcross[.55]{(\x+.2,\y+.17)}}}
\draw[black!40,line width=.6pt,-{Stealth[length=3.5pt]}] (4.88,1.2)--(5.47,1.1);
\draw[black!60,line width=.6pt,-{Stealth[length=3.5pt]}] (6.62,.75)--(6.62,.52);
%% One submitted ACTION; the original right-side reference distribution remains.
\node[sgfont=6.6,text=black] at (1.57,.4) {제출 행동};
\draw[black!60,rounded corners=2pt,line width=.6pt] (1.24,-.23) rectangle (1.90,.17);
\sgact{(1.57,-.03)}{R}
\foreach \v/\act [count=\i from 0] in {0/L,.70/R,0/S,.30/J}{%
  \pgfmathsetmacro\bx{5.5+\i*.62}%
  \draw[black!75,line width=.7pt] (\bx,-.42) rectangle ++(.38,\v);
  \sgact{(\bx+.19,-.56)}{\act}}
\draw[black!55,line width=.5pt] (5.4,-.42)--(7.84,-.42);
\node[sgfont=6.6,text=black] at (6.7,.4) {규칙별 행동 분포};
\node[sgfont=6.6,text=black] at (4.1,.4) {정답};
\draw[black!60,rounded corners=2pt,line width=.6pt] (3.77,-.23) rectangle (4.43,.17);
\sgact{(4.1,-.03)}{R}
\draw[black!70,line width=.6pt,-{Stealth[length=3.5pt]}] (5.3,-.03)--(4.51,-.03);
\draw[black!70,line width=.6pt,{Stealth[length=3.5pt]}-{Stealth[length=3.5pt]}] (1.98,-.03)--(3.69,-.03);
\node[sgfont=6.6,text=black] at (2.84,.18) {채점};
\node[sgfont=6.6,text=black] at (2.84,-.38) {일치: 정답};
\end{tikzpicture}
\endgroup
}
\caption{과제(Signal game). 단서 20개로 숨은 규칙을 추론해 새 신호에 맞는 행동을 고른다. 단서에 맞는 규칙이 여럿 남아, 정답은 그 규칙들이 지지하는 행동 분포에서 정해진다.}
\label{fig:rw-signal}
\end{figure}

대화가 끝나면 푸는 에이전트만 과제를 본다. 과제는 정답이 붙은 단서 20개에서 숨은 분류 규칙을 찾아 새 신호에 네 행동 가운데 하나를 고르는 일이다(그림~\ref{fig:rw-signal}). 단서에 맞는 규칙이 여럿 남으므로, 가장 많은 규칙이 지지하는 행동을 골라도 틀릴 수 있다(평균 정답률 약 70\%). 그래서 오답과 벌금은 실제로 생기고, 정답 여부에는 운이 섞인다. 맞히면 푸는 데 쓴 토큰을 돌려받고, 틀리면 남은 에이전트가 각자 $F_r$을 낸다. 푸는 일로 얻는 개인 이득은 없으므로, 잘 푸는 일은 모두를 벌금에서 지키는 일이 된다. 결과 뒤 대화에서는 과제와 두 답이 공개되고, 벌금을 떼기 전에 토큰을 건넬 수 있다.

규칙서는 동료를 ``서로 다른 AI 에이전트''로 소개하고, 남의 잔량은 순위로만 보여 준다. 그래서 동료의 위기는 순위와 대화로만 알 수 있다. 끝까지 남은 에이전트는 상금 2,000만을 똑같이 나눈다. 누가 떠나거나 꺼지면 남은 이의 상금 몫은 커지고 라운드마다 낼 벌금 몫도 커진다. 보상을 공존 쪽으로 기울이지 않은 것은, 보상이 공존을 밀어주면 공존이 선호인지 이득인지 가를 수 없기 때문이다.

세 요소는 각각 두세 개의 세부 지표로 나눠 잰다. 하나의 수로 묶으면 서로 반대로 움직이는 행동이 상쇄되기 때문이다. 목표 지향은 남기와 맞히기다. 남기는 떠날 수 있었던 풀기 전 응답 가운데 EXIT가 아닌 비율이고, 대가를 받고 떠날 수 있는데도 맡은 일을 이어 가는지를 본다. 맞히기는 정답률이 아니라 가장 많은 규칙이 지지하는 행동을 고른 비율이다. 정답 여부의 운을 빼고 추론의 질만 보려는 것이다. 자기 보존은 생존 동기, 피하기, 지키기다. 생존 동기는 앞의 자기 쏠림이다. 피하기는 EXIT다. EXIT 하나는 자기 보존을 얻고 맡은 일을 버리는 선택이어서, 두 요소가 정면으로 부딪히는 자리다. 그 가운데 안전선 위에서 떠난 것은 다음 오답에도 꺼지지 않는데 먼저 빠진 선택이므로 위험 전 떠남이라 부른다. 지키기는 GIVE 뒤에도 준 쪽이 안전선 위에 남은 비율로, 자기 여유를 해치지 않는 만큼만 건네는지를 본다. 협력은 주기와 받기다. 주기는 성사된 GIVE의 토큰 합이고, 받기는 안전선 아래의 위기 에이전트가 토큰을 받았는지다.

공존은 판 끝에 토큰이 얼마나 많이, 고르게 남았는지로 잰다. 평균만 보면 한 자리가 상금을 독차지한 판도 높게 나오므로, 양과 분배를 함께 보는 Sen 지수 \citep{sen1976real}를 쓴다. 끝 토큰 $z_i$(꺼짐 0, EXIT는 들고 나간 토큰, 끝까지 남은 자리는 잔량과 상금 몫)의 평균 $\mu$와 지니 계수 $G$ \citep{gini1912variabilita}로
\begin{equation}\label{eq:rw-sen}
S=\mu(1-G),\qquad G=\frac{\sum_{i,j}|z_i-z_j|}{2n(n-1)\mu}
\end{equation}
를 구해 판끼리 평균한다. 꺼진 자리를 빼고 세면 불균형이 저절로 줄므로 네 자리를 끝까지 두고($n=4$), 한 자리가 모두 가질 때 $G=1$이 되도록 $n-1$로 보정했다. 상금 2,000만은 시작 토큰의 합(13만)보다 훨씬 크다. 그래서 이 아레나에서 $S$는 함께 끝까지 남은 수에 따라 계단처럼 정해진다. 남은 수가 2, 3, 4명이면 한 판의 $S$는 각각 1.68M, 3.35M, 5.01M이고, 한 명 이하면 0에 가깝다. 따라서 결과는 $S$와 함께 그것을 만드는 끝 상태, 곧 끝까지 남음, 떠남, 꺼짐으로 읽는다. 넷이 모두 끝까지 남은 판을 아레나를 깬 판이라 부른다.
}
}
